% 第七章：工程问题与解决过程

\section{工程问题与解决过程} \label{sec:difficulties}

本章按"现象—根因—处理—证据"记录七个真实问题。它们的共同点是：都能靠"跑一遍、对一遍数"发现，靠"读一遍代码"定性。更重要的是每个问题最后都留下了一条门禁或一个开关——工程上的收获不是"踩过坑"，而是此后同类问题会被自动拦住。

\subsection{十个合并冲突：路线选择}

\textbf{现象}：\texttt{kernel/sched/yat\_casched.c} 遗留 10 处冲突标记，两套实现（初赛版与决赛版）交织，直接编译会报语法错。

\textbf{根因}：并行开发时同一文件被两条演进线修改，合并没有按语义收敛，而是按行保留了两边。

\textbf{处理}：取"决赛 AJLR 版"为唯一事实，逐处判断哪些初赛函数仍被调用、哪些只是历史残留；删除无调用者的死代码，保留被 \texttt{sched\_class} 或测试引用的接口。收敛后 \texttt{make symbol-check} 通过，头文件 29 个声明全部有定义。

\textbf{证据}：\texttt{docs/CONFLICT\_ANALYSIS.md} §7--§8 记录逐处决策。

\subsection{早期仿真数据的口径问题（已更正）}

\textbf{现象}：项目前期对外宣称的 Makespan 改进 11.66\%、命中率提升 30.21 个百分点、胜率 96\%，在复核后无法复现；用同一份代码重跑，数字变成 $-0.954\%$，换一个入口再跑又变成 $-3.524\%$。

\textbf{根因}有三层，一层比一层靠近真问题：

\begin{enumerate}
  \item \textbf{两套模型口径}。旧实验里 CacheAware 用自估函数把执行时间打折（最高减免 50\%），WFD 用原始值；命中率不由缓存模型给出，而是分析器按算法名套两套手写公式，其中 WFD 那套还被截断到 $[0.30,0.70]$。所谓"更快、更准"是口径差异的产物。
  \item \textbf{基线状态泄漏}。\texttt{WFD} 的处理器负载表是 \texttt{static}，只在每个种子结束时清理一次，于是第 $i$ 个用例的基线继承了前 $i-1$ 个用例累积的负载偏置，而挑战者每个用例都复位。同一批用例，这个偏差稳定地把结果抬高约 2.5 个百分点。
  \item \textbf{决策变量缺项}。旧启发式的评分为五维加权，五项全部来自任务的静态属性（敏感度、权重、随机亲和性、自维护的"缓存利用率"），没有任何一项依赖"这个 DAG 上次在哪个核跑过"；而缓存模型的命中判据只有一条——最近 5 次访问中是否有同一 \texttt{dagID}。两边正交，任何权重调参都无法把局部性变成收益。
\end{enumerate}

\textbf{处理}：把仿真器整体移出证据链。它的执行时间与命中判定都由自建模型给出，不含抢占、迁移、TLB 与真实 LLC 竞争，即使口径全部修正，也只是"理想模型下的行为"，不能支撑"在真实系统上有效"的结论。正文只保留第\ref{sec:test}节的一小段交代，数字与复现命令留在 \texttt{docs/validation/}。

\textbf{证据}：\texttt{docs/validation/2026-11-13/SUMMARY.md} 记录三套口径（11.66\% / $-0.954\%$ / $-3.524\%$）的逐种子对照与逐字节复现结果；仓库同时保留 \texttt{--legacy-metrics} 与 \texttt{--legacy-wfd-state} 两个开关，使历史数字随时可对齐——保留它们的目的不是"留个说法"，而是让"哪一版对、差在哪"可被第三方独立验证。

\subsection{QEMU 的两个参数坑}

\textbf{现象 1}：\texttt{-kernel vmlinux} 直接报 \texttt{Error loading uncompressed kernel without PVH ELF Note}，不启动。

\textbf{根因}：QEMU 对 x86 ELF 内核要求 PVH note，本项目的 vmlinux 没有；\texttt{bzImage} 走的是另一条加载路径。

\textbf{处理}：一律用 \texttt{arch/x86/boot/bzImage}。

\textbf{现象 2}：加了 \texttt{acpi=off} 后 guest 只有 1 个 CPU（\texttt{nr\_cpu\_ids=1}），"多核调度"实验失去意义。

\textbf{根因}：\texttt{-kernel bzImage} 绕过 BIOS，AP 枚举只能靠 ACPI MADT；关掉 ACPI 就没有 AP。

\textbf{处理}：保留 ACPI，不加 \texttt{acpi=off}/\texttt{noapic}/\texttt{nolapic}，实测 4 核正常起来且 \texttt{poweroff} 能真正退出。这两条写进脚本头注释，避免后人"顺手优化"回去。

\subsection{锁反序：lockdep 报出来的环}

\textbf{现象}：镜像 \#18 的串口日志在基准跑完后立即出现 \texttt{WARNING: possible irq lock inversion dependency detected}，涉及 \texttt{pool->lock} 与 \texttt{rq->\_\_lock}。

\textbf{根因}：\texttt{pool->lock} 在 balance 路径里总是关中断获取（HARDIRQ-safe），而 \texttt{task\_dead} 在开中断上下文里获取它，构成"任务持池锁 $\to$ 中断 $\to$ 取 rq 锁 $\to$ 需要池锁"的环。此前只把 debugfs 的 \texttt{show()} 改成 \texttt{irqsave}，属于堵旁路而非断环。

\textbf{处理}：把 \texttt{task\_dead} 里的整段销毁放进 \texttt{local\_irq\_save()}/\texttt{local\_irq\_restore()}。临界区内只有链表操作与 \texttt{kfree}，没有睡眠式操作。

\textbf{证据}：修复后重建（镜像 \#19，SHA256 \texttt{077638d4\ldots}），复跑两组对照，串口日志中同类告警计数为 0；修复并入发布补丁 \texttt{0001}，用"修复前文件 + 补丁 == 当前文件"做过往返验证。

\subsection{CRP 归一化窗口：一处注释错误，牵出一个更严重的问题}

\textbf{现象}：\texttt{yat\_get\_crp\_ratio()} 里 \texttt{max\_recency = 1000000}，注释写"1 秒，根据实际情况调整"。\texttt{ktime} 的单位是纳秒，$10^6$\,ns 是 1\,ms，不是 1\,s。

\textbf{修复}：把常量提为 \texttt{YAT\_CRP\_MAX\_RECENCY\_NS}，注释改成 1\,ms 并标注"未经标定"；新增两个计数器 \texttt{crp\_used} 与 \texttt{crp\_sat}，把"比率算了几次、其中几次撞上窗口"变成可读数。

\textbf{更严重的发现}：有了计数器之后，在主口径上连跑三轮，饱和率分别是 96.6\%、97.4\%、5.3\%。也就是说，前两轮里几乎每个候选核的 $\mathrm{benefit}$ 都恒为 0，选核实际退化成按 WCET 排序——第\ref{sec:test}节那两条成本数（$1.81$ 与 $1.72$）描述的是"任务经过这个类"的路径开销，而不是"缓存决策"的代价。这与空类标定（把记账与均衡一起摘掉后比值不降反升）指向同一结论。

原因不在那个常数，而在结构：$\mathrm{recency}=L_1+2L_2$，而 $L_2$ 是\emph{所有核}自本任务上次运行以来的执行时间之和。8 个任务跑在 4 个核上时，任务离开核不到 1\,ms，各核累加值就超过窗口；饱和率出现 5.3\% 的那一轮，只是唤醒节奏更碎、离开核更短而已。换言之，CRP 的分辨力取决于负载相位，而不是模型是否"正确"。

\textbf{处理}：此刻仍不改窗口数值——没有标定数据就调参，正是本章第二节批评的做法。改为把"饱和率"登记为可自动检查的判据（$\texttt{crp\_sat}/\texttt{crp\_used}<20\%$），并把结构性的修法（$L_2$ 按在线核数归一）列为路线 D 的第一步。

\subsection{文档漂移}

\textbf{现象}：同一组数字在 README、\texttt{EXPERIMENT\_REPORT.md}、\texttt{ROADMAP.md}、\texttt{index.html} 里出现过 11.66/30.21/96\%、0.39/0.82/56\%、$-0.954$/0.374/48.6\% 三种口径，区间也有 $\pm1.745$ 与 $\pm1.23$ 两版。

\textbf{处理}：以 CSV 重算值为唯一事实，把三种口径的适用条件写进各自文档；\texttt{make docs-check} 校验链接、\texttt{make check} 校验冲突标记；一致性靠单一引用点（\texttt{docs/validation/2026-11-13/SUMMARY.md}）加各文档顶部的时效声明，而不是靠人工同步。

\subsection{两处工具链故障，与它们留下的门禁}

\textbf{现象 1}：补丁包 	exttt{0005} 在发布前被发现无法应用——	exttt{git apply} 报 	exttt{corrupt patch at line 3276}，最后一段 hunk 被截断。而当时 	exttt{make patch-check} 一直是"通过"的：它只校验 	exttt{MANIFEST.sha256} 与 hunk 数量，不校验能否应用。

\textbf{根因}：补丁是从工作树一次性导出的，生成脚本在写最后一个文件时被中断，文件少了一个尾行；校验项恰好不覆盖"可应用性"。

\textbf{处理}：改为按当前树逐文件重新生成，然后在临时目录里应用整个 series 并逐文件比对，确认 21 个文件与工作树一致。补丁往返（应用 $	o$ 比对）现在是发布前的固定步骤；同时明确一点：	exttt{sha256} 只能证明"文件没被改"，不能证明"补丁是对的"。

\textbf{现象 2}：同时起两个对比脚本，两边共用 	exttt{results/} 下的 initramfs 与串口日志，一方会读到另一方的日志，表现为"未在日志中找到 	exttt{YATCSV} 行"的假失败，严重时两个 QEMU 同时段错误。

\textbf{根因}：测量脚本默认"一次只跑一个实例"，没有任何互斥；而实际使用中并行跑是常态。

\textbf{处理}：脚本加 	exttt{flock}，第二个实例直接退出并提示锁路径。这条改动的意义不在脚本本身，而在于它把"测量必须独占"从口头约定变成机器约束——一份会被并行污染的实测数据，比没有数据更危险。

\textbf{证据}：两项修复后，15 次重复测量（3 配置 $	imes$ 5 次）连续跑完、无假失败、无段错误，结果与 	exttt{kernel/repeat-routeA/} 归档一致。
